% MAS1004 Assignment 2: long report.
%
% You fill this in with your agent, in your repository. AGENTS.md tells the
% agent how. Red text in the PDF is what still has to be replaced.
%
% - Pictures go in figures/. Copy them from results/, which is not in git.
% - Outputs go in outputs/: copy experiments.csv, LOG.md and the outputs of
%   check_submission.py and curate.py count here. \outputfile{outputs/...}
%   prints a file exactly as it is. Every number in this report has to come
%   from such an output.
%
% To make the PDF: compress this folder into a zip, and on Overleaf choose
% New Project, Upload Project. Set Menu, Compiler to XeLaTeX, and Recompile.

\documentclass[11pt]{article}
\usepackage{mas1004}

\reporttitle{Assignment 2: Long report}
\studentname{Your name}
\studentid{Your student ID}

\begin{document}
\makeheader

\fillin{Your leaderboard nickname, and the tags of your last Track A and Track B releases.}

% The sections below are a suggestion. Change, merge or add sections as you
% like; README.md, "The long report", says what a reader should be able to see.

% ---------------------------------------------------------------------------
\section*{Track A}

\begin{center}\scriptsize
\begin{tabular}{p{0.17\linewidth}p{0.33\linewidth}rrrrrr}
\toprule
tag & what changed & validation & seoul & busan & tokyo & osaka & taipei \\
\midrule
baseline (seed 0) & nothing, the starter recipe, trained with seeds 0 and 1 & 0.4961 & 0.5774 & 0.3763 & 0.6259 & 0.343 & 0.5577 \\
baseline (seed 1) & nothing, the starter recipe, trained with seeds 0 and 1 & 0.4922 & 0.5622 & 0.3871 & 0.606 & 0.3736 & 0.5319 \\
more\_often\_outside & in sample\_weights, photos whose file name starts with osv5m\_ or youtube\_ are drawn twice as often as the others (commons\_); no photo removed & 0.4876 & 0.5651 & 0.3691 & 0.5988 & 0.3802 & 0.5247 \\
balanced & back to the starter recipe (more\_often\_outside undone), then in sample\_weights every city drawn equally often: each photo gets one over the number of photos of its city & 0.5117 & 0.5113 & 0.4424 & 0.5256 & 0.4511 & 0.6282 \\
balanced\_zoom & on top of balanced, RandomResizedCrop scale from (0.35, 1.0) to (0.1, 1.0) & 0.5089 & 0.4723 & 0.4755 & 0.5125 & 0.4492 & 0.6349 \\
balanced\_brightness & back to balanced (balanced\_zoom undone), then ColorJitter brightness from 0.2 to 0.5 (each training photo made between 0.5 and 1.5 times as bright); contrast and saturation stay 0.2 & 0.5064 & 0.4846 & 0.4511 & 0.5014 & 0.4629 & 0.6322 \\
balanced\_rule\_drop & recipe back to balanced (balanced\_brightness undone); 61 training photos dropped by my Step 7 rule (busan 11, osaka 10, seoul 14, taipei 11, tokyo 15), reason "closeup or focus on one thing" & 0.4952 & 0.5027 & 0.4302 & 0.5259 & 0.4272 & 0.5903 \\
balanced\_smaller\_lr & on top of balanced (all training photos again), peak learning rate from 3e-4 to 2e-4 (lr in make\_optimizer and max\_lr in make\_scheduler); schedule, warm-up and weight decay unchanged & 0.5065 & 0.5224 & 0.4525 & 0.513 & 0.4413 & 0.6032 \\
balanced\_zoom\_small\_lr & balanced with two earlier changes combined: RandomResizedCrop scale (0.1, 1.0) as in balanced\_zoom, and peak learning rate 2e-4 as in balanced\_smaller\_lr; nothing else & 0.5138 & 0.4932 & 0.4978 & 0.5022 & 0.437 & 0.6389 \\
balanced\_zoom\_lr\_outside\_tokyo & on top of balanced\_zoom\_small\_lr, within Tokyo only, photos whose file name starts with osv5m\_ or youtube\_ are drawn twice as often as Tokyo's other photos; every city still gets the same share & 0.5101 & 0.5055 & 0.4748 & 0.4762 & 0.4718 & 0.6219 \\
\bottomrule
\end{tabular}
\end{center}

Each run, with my reason and what I think the result means, in my own words from the lab notebook:
\begin{itemize}\small
\item \textbf{baseline} (validation 0.4961). Why: ``I trained the starter recipe with seed 0 and 1 and the tag baseline as the step 2 said. After doing step 1 and looked at the photos I noted what could tell the cities apart, but decided to first train the starter recipe and later make changes. In step 3 I can look at the mistake ofthe starter recipe, and have a better idea for the changes.'' What it means: ``From what I understood, a seed changes the random starting point, so it's telling me there is a slight difference based on where the starting point it, but i'm not sure is telling much else since it is a very small change.''
\item \textbf{more\_often\_outside} (validation 0.4876). Why: ``For the first change I think it's good to start from what I noticed in the mistake sheets. The first element was that it was outside photos, second was photos with no closeups. I would do one change at a time. Since the model got wrong many outside photos it could be useful to train it on a set of data where the outside images are many and appear more times. I think I wouldn't delete the inside photos but it would be good to make the outside photos appear more often during the training. / yes, let's use the name of the file since doing it by checking each photo is not possible. This way it should be mostly outside pictures. So let's do it for both youtube and osv5m, twice as often for now.'' What it means: ``I think this results tells me that the outside photo are perhaps a problem. i thought that training the model on more photos of the outside would be useful but the validation was lower than the starting recipe, only Osaka had a higher validation.''
\item \textbf{balanced} (validation 0.5117). Why: ``let's try to make the model train on the same amount of photo for each city. It could be useful for the model to train equally on each city photos, and it could be a good point for then observing changes and what mistakes it makes in that condition'' What it means: ``This results tells me that it's good for the model to train on equal variety of photos from each city, even if it means that Seoul and Tokyo got lower score.''
\item \textbf{balanced\_zoom} (validation 0.5089). Why: ``I think for the next change it would be good to work on balanced to maintain the improvement achieved with it and see the effects of next changes when cities are equally drawn. I think I want to try the zoom in the photos since some mistake photos in balanced are still from afar so I would still want to try since it was one of the first thing I noticed at the beginning.'' What it means: ``The results tell me that adding the zoom to the balanced didn't improve the model overall, although there are some improvements for Busan and Taipei, but Seoul and Tokyo went down a bit. So the zoom was useful for some city but not for others, meaning it's not an overall change that works for all the cities. It could be kept into consideration to combine in a second moment with other changes.''
\item \textbf{balanced\_brightness} (validation 0.5064). Why: ``for the next change I would start from balanced again, in order to check the change without zoom. I think at this point I should address another element I noticed in the mistake sheets, which is that some photos in the mistake are darker. So it could be interesting to make some photos lighter for the training, perhaps by making the light vary like for the zoom / let's use both, darker and lighter 0.5, - 1.5 , since in the mistakes for some cities the photos were darker, while in other not.'' What it means: ``this results tell me that change in brightness were less useful than zoom, but it still was slighly good for Busan Osaka and Taipei. For Osaka it was more effective than zoom. It shows a differnet situation for Seoul and Tokyo, which makes me think I should consider these two city in similar way since both went down and up together with changes.''
\item \textbf{balanced\_rule\_drop} (validation 0.4952). Why: ``my Step 7 rule above'' What it means: ``the validation on this training was lower than the normal balanced so it makes me think that perhaps what I decided to drop and my rule was not the best choice.''
\item \textbf{balanced\_smaller\_lr} (validation 0.5065). Why: ``I think I want to try to change the learning rate on balanced. I remember from assignment 1 that a larger learning rate was not necessarily better so in this case I would make it a bit smaller. / I wouldn't go too much lower so 2e-4 for now'' What it means: ``the results of this training are lower than the balanced one, it was good for Seoul but not for Taipei. It seems like the changes I made after balanced, were less effective or were only effective for some cities but not for the overall validation.''
\item \textbf{balanced\_zoom\_small\_lr} (validation 0.5138). Why: ``At this point I think it could be interesting to try and combine previous changes since after balanced, the other changes were positive for some cities and negative for others. Perhaps combining the change in zoom (which previously was positive for Busan and Taipei especially) with smaller lr (positive for seoul), could create an overall better validation than balanced by itself. The risk is that the two changes cancel each other and the validation ends up still lower, but in that case it would tell me that it's not the best way to combine.'' What it means: ``combining zoom and smaller lr was useful, especially for Busan and Taipei. I expected Seoul to have a higher score since with the learning rate it increased. I also thought tokyo would stay around the scores of zoom and small lr, but didn't, meaning I still have to find what makes Tokyo score go up.''
\item \textbf{balanced\_zoom\_lr\_outside\_tokyo} (validation 0.5101). Why: ``I'm thinking of making the more often outside for tokyo (since in that change previously tokyo was higher than in balanced) / Then, I want to try only more outside for tokyo, since I'm more interested now to see how the more outside work on the balanced. / use the same twice as often for the outside for tokyo'' What it means: ``The change in showing more outside for tokyo didn't work the same as when it was applied to the baseline, meaning that balanced made a difference, and I still need to find what on balanced makes Tokyo go up.''
\end{itemize}

\fillin{Anything you want to add about the Track A runs as a whole.}

\begin{center}\footnotesize
\begin{tabular}{lllrrrrrr}
\toprule
release & title & & score & seoul & busan & tokyo & osaka & taipei \\
\midrule
track-a-1 & starter recipe & leaderboard & 0.400 (0.395--0.403) & 0.62 & 0.19 & 0.70 & 0.25 & 0.23 \\
 & & validation (grader) & 0.497 & & & & & \\
track-a-2 & balanced & leaderboard & 0.410 (0.402--0.419) & 0.57 & 0.26 & 0.61 & 0.33 & 0.28 \\
 & & validation (grader) & 0.502 & & & & & \\
track-a-3 & balanced zoom+lr & leaderboard & 0.403 (0.397--0.414) & 0.55 & 0.28 & 0.61 & 0.30 & 0.28 \\
 & & validation (grader) & 0.509 & & & & & \\
\bottomrule
\end{tabular}
\end{center}
The leaderboard score is the mean of the three models the grader trains (seeds 0, 1 and 2), with the lowest and highest in brackets; the grader's validation is their mean on \texttt{data/starter/val}.

\fillin{Your final recipe (the last release before the deadline), and why you chose it.}

% ---------------------------------------------------------------------------
\section*{Decisions you made from photos}

\picturefile{figures/suspects_balanced_seed0_tokyo_00.png}

I looked at the suspects sheets of balanced (\texttt{suspects\_balanced\_seed0\_<city>\_00}): for each city, the 48 training photos, out of every fourth one, that the balanced model gave the lowest probability for their own city.

What I saw, from my lab notebook:
\begin{quote}\small
Suspect balanced Busan: I notice photos that are closeup or focus on object or things such as food, vase, flowers, books and vehicles.

Suspect balanced Osaka: similarly photos closeup or focus of animal, food

Suspect balanced Seoul: closeups on objects, food. but also outside buildings and people

Suspect balanced Taipei: here similarly photos of objects, animal, food,

Suspect balanced Tokyo: I notice photos of food, cars, people
\end{quote}

The rule I wrote before dropping any photo, from my lab notebook:
\begin{quote}\small
I think overall I notice that there als also many outside photos, but I also think that ouside photos usually have a lot of elements into them that can be useful for the model, so I wouldn't drop them. I think finding one specific only rule is not easy as the images a various, so I think the rule I want to make follows what I noticed, meaning photos that are closeups/focus on only one object or thing, that doesn't give much other element in the photo.
\end{quote}

After dropping the 61 photos, \texttt{curate.py count} printed:
\outputfile{outputs/curate_count.txt}

Validation before (balanced): 0.5117; after (balanced\_rule\_drop): 0.4952. Afterwards the 61 photos were put back.

What I think it means, from my lab notebook:
\begin{quote}\small
the validation on this training was lower than the normal balanced so it makes me think that perhaps what I decided to drop and my rule was not the best choice.
\end{quote}

My requests to the agent for these decisions, word for word:
\begin{itemize}\small
\item Drop 3, 16, 17, 20, 26, 28, 30, 32, 34, 44, 48, from suspects\_balanced\_seed0\_busan\_00.png because “closeup or focus on one thing”
\item Drop 5, 6, 9, 10, 11, 16, 20, 34, 36, 48 from suspects\_balanced\_seed0\_osaka\_00.png  because “closeup or focus on one thing”
\item Drop 1, 3,4, 7, 8, 18, 20, 22, 23, 26, 30, 33, 41, 43, from suspects\_balanced\_seed0\_seoul\_00.png because “closeup or focus on one thing”
\item Drop 4, 5, 12, 15, 17, 26, 30, 31, 33, 36, 42, from suspects\_balanced\_seed0\_taipei\_00.png because “closeup or focus on one thing”
\item Drop 6, 9, 10, 14, 17, 18, 22, 23, 25, 31, 36, 39, 41, 42, 43, from suspects\_balanced\_seed0\_tokyo\_00.png because “closeup or focus on one thing”
\end{itemize}

% ---------------------------------------------------------------------------
\section*{Track B}

\begin{center}\scriptsize
\begin{tabular}{p{0.17\linewidth}p{0.33\linewidth}rrrrrr}
\toprule
tag & what changed & validation & seoul & busan & tokyo & osaka & taipei \\
\midrule
track\_b1 & nothing, train\_b.py as it is (ResNet18 at 224, 4 epochs, starter training photos) & 0.4933 & 0.593 & 0.3835 & 0.6127 & 0.3651 & 0.5123 \\
track\_b\_balanced & train\_b.py with -{}-balance (every city drawn equally often); everything else as track\_b1 (ResNet18 at 224, 4 epochs, starter training photos) & 0.4976 & 0.5314 & 0.4403 & 0.4878 & 0.4521 & 0.5765 \\
track\_b\_balanced\_smaller\_lr & on top of track\_b\_balanced, peak learning rate from 3e-4 to 1.5e-4 (-{}-lr 1.5e-4); everything else the same (ResNet18 at 224, 4 epochs, -{}-balance) & 0.4955 & 0.5047 & 0.4489 & 0.5133 & 0.4445 & 0.5662 \\
\bottomrule
\end{tabular}
\end{center}

\begin{itemize}\small
\item \textbf{track\_b1} (validation 0.4933). Why: ``this run is the requested from step 8, train track b as it is.'' What it means: ``It's slighltly different from the track A baseline, but still somewhat similar, where Tokyo has highest score, seoul is slightly higher than the baseline, and the same is for Busan, and Osaka.''
\item \textbf{track\_b\_balanced} (validation 0.4976). Why: ``I want the first change for track B to follow what worked for track A, meaning balance, so that each city photos are drawn equally. The reason is that as a first step I think it's good to go with what worked in the other track as well, before trying other changes.'' What it means: ``the results show that as for track A balanced improved the model, especially for Busan Osaka and Taipei.''
\item \textbf{track\_b\_balanced\_smaller\_lr} (validation 0.4955). Why: ``fro track B let's try to change the learning rate as for track A but this time instead of two third (2e-4) I want to try half (1.5e-4). Since smaller lr worked for track A for seoul, and busan, this could be interesting to see if in this case it's too small or not, and can be useful overall.'' What it means: ``Compared to track A, a smaller learning rate (in this case half) for track B didn't bring the validation down significantly, which leaves open whether using a smaller learning rate, like in track A, gives a different result, and could be useful for the model.''
\end{itemize}

No extra photos: every Track B model was trained on the starter training photos only.

\begin{center}\footnotesize
\begin{tabular}{llrrrrrr}
\toprule
release & & score & seoul & busan & tokyo & osaka & taipei \\
\midrule
track-b-1 (track\_b1) & leaderboard & 0.389 & 0.60 & 0.16 & 0.68 & 0.23 & 0.28 \\
 & validation & 0.4933 & 0.593 & 0.3835 & 0.6127 & 0.3651 & 0.5123 \\
track-b-2 (track\_b\_balanced) & leaderboard & 0.405 & 0.63 & 0.24 & 0.61 & 0.30 & 0.24 \\
 & validation & 0.4976 & 0.5314 & 0.4403 & 0.4878 & 0.4521 & 0.5765 \\
\bottomrule
\end{tabular}
\end{center}

\fillin{Why you think validation and leaderboard disagree (your first thoughts are in the lab notebook, under Leaderboard).}

\subsection*{Your submission, as the grader sees it}
\texttt{check\_submission.py --no-val} on the files of release track-b-2:
\outputfile{outputs/check_submission.txt}

% ---------------------------------------------------------------------------
\section*{What made the difference}

\fillin{Which change helped most and which least, in both tracks, with your
numbers.}

\fillin{What you would do next with one more week.}

% ---------------------------------------------------------------------------
\section*{All runs}
\outputfile{outputs/experiments.csv}

\section*{Lab notebook}
\outputfile{outputs/LOG.md}

\end{document}
